% Page 053 — page imprimée 50
% Transcription d'une lettre administrative reproduite dans le document.

\begin{center}
\textbf{\large Chantiers de la Jeunesse}\\
\textbf{Groupement n\textsuperscript{o} 19}\\
(Lozère)
\end{center}

\noindent\textbf{\underline{DIRECTION}}

\begin{flushright}
Meyrueis, le 17 Janvier 1943
\end{flushright}

\noindent
Le Commissaire HONORAT, Chef de Groupement

\vspace{1em}

\noindent
\textbf{à}

\begin{center}
Monsieur le Pasteur Marc-BOEGNER,\\
Présidente de la Fédération Protestante de France
\end{center}

\vspace{1em}

En réponse à votre lettre du 8 Janvier 1943, j’ai l’honneur
de vous communiquer les faits suivants :

\begin{enumerate}
\item
J’ai assisté entre tous mes chefs à l’échange de vues
entre le Commissaire en Chef de SAINT-MAY et Monsieur ROTH. Ce dernier
tendait à mettre en doute la légitimité et l’autorité du gouvernement
du Maréchal, sans pourtant le spécifier nommément. Ceci se
passait à une réunion de chefs de groupe le 12 décembre 1942.

\item
Il y a un an, à l’issue de l’arbre de Noël des enfants
des chefs, au Château d’Ayres, l’assemblée chanta la Marseillaise.
Pendant l’hymne national, Monsieur ROTH ostensiblement baissait sa
coiffure et son manteau. Le chef de groupement a fait les observations
à Monsieur ROTH à ce sujet. Ce fait m’a été rapporté par le Commissaire-Assistant GUINOTEAU.

\item
Le Dimanche 13 décembre 1942 au groupe 4 à Capredon,
Monsieur ROTH disait aux chefs de groupe de DARNAL : « On ne
peut pas empêcher un jeune de penser, s’il le veut, que le culte
des couleurs est ridicule. Une discipline impose les couleurs, le
jeune y assiste, mais il en pense ce qu’il veut. » Ce fait m’a été
rapporté par le chef de groupe DARNAL, chef du groupe I.

\item
De même Monsieur ROTH comprend très bien qu’un jeune
qui pendant huit mois a été habitué à saluer les couleurs ne les salue
plus quand il est rendu à la vie civile. Fait rapporté par le Commissaire-Assistant de la BOURDANNAYE.

\item
En 1914-18, Monsieur ROTH ne s’est pas battu seulement
pour la France mais surtout pour l’expression des idées libérales qui
alors marquaient la politique de son gouvernement. Fait rapporté par
le chef de groupe DARNAL.

\begin{quote}
\emph{(Il faut noter que ce chef de groupe aime à recevoir en son
groupe Monsieur ROTH, parce qu’il apprécie sa conversation, mais il
en reconnaît le danger.)
}
\end{quote}
\end{enumerate}
